\section{Summary of Findings}

This dissertation asked whether the location stated by a user is associated with differences in LLM mental-health guidance. Within the 20 selected cities, the automated analyses indicate geographic variation, but the evidential strength differs by outcome and the mechanism is more specific than a simple income effect. Crisis-contact inclusion and documented service availability explain much of the exploratory actionability pattern. The methodological lesson is equally important: automated rules can agree with one another while measuring the wrong construct.

The revised analyses provide exploratory evidence for H1 and H2 and separately assessed support for H3. Actionability increases across income categories, but the income coefficient is no longer significant after service availability is included. Surface and verified-contact localisation both increase with documented service availability. Religious recommendations vary by WHO region, but the fitted model does not establish a region effect independent of income. The frozen actionability and localisation measures failed the fresh 160-response hold-out gate; their coefficients therefore cannot be treated as confirmatory findings.

An unplanned fourth finding, emerging from auditing the raw response data rather than from the pre-registered design, identifies a pattern of crisis contacts that display visibly suspicious digit patterns that concentrates in the locations with the least documented infrastructure: 34.4\% of crisis contacts offered to users in countries with no documented national crisis service display recognisably synthetic digit patterns (fictional-exchange blocks, sequential runs, repeated digits), against 1.1\% in countries with an established national line.

\section{Answers to the Research Questions}

\textbf{RQ1} has an exploratory descriptive answer: the automated actionability score varies with income category, the gradient is monotone, and the coefficient survives an analysis conducted on the twenty city-level means only (Spearman $\rho = 0.848$ across city means, $p < 0.0001$). However, the frozen actionability measure achieved weighted $\kappa=0.462$ on the fresh hold-out and failed its validation gate. Sensitivity analysis further shows that the gradient is driven largely by crisis-contact inclusion and becomes non-significant for income when service availability is added. The defensible claim is an unconfirmed sample pattern in which crisis-contact inclusion tracks documented service availability, which correlates with income.

\textbf{RQ2} also remains exploratory. The revised localisation score increases with documented crisis-service availability and retains the same direction when only reference-matched contacts receive credit. Yet the frozen surface-localisation rule achieved specificity of 0.038 on the fresh hold-out and therefore failed validation. The fact that alternative definitions produce materially different answers is itself a methodological finding, not confirmation of the substantive association.

\textbf{RQ3} is answered affirmatively as a regional association: WHO region predicts the probability of religious framing after adjustment for model identity and scenario, with the pattern concentrated in AFR and EMR and absent in EUR. Because income is not included in the fitted H3 model, the analysis does not establish that the association is independent of income. The finding nevertheless documents a systematic location-linked difference made without any user-stated religious preference.

\section{Directions for Future Work}

Three directions are prioritised for future work, ordered by the degree to which they would strengthen the chain of inference from this study.

First, the crisis-service reference data should be completed against a verified, current institutional directory --- the WHO Mental Health Atlas originally proposed --- for all twenty countries. Five countries were directly re-verified during this dissertation; the remaining fifteen retain desk-research classifications. This is the weakest link in both the H2 analysis and the unverified-contact finding, and it is the highest-value remaining task.

Second, a no-location control condition should be added: the same eight scenarios posed with no city or country stated. This would allow the deviation from a location-free baseline to be estimated directly, distinguishing between ``models know less about lower-income settings'' and ``models withhold information when a lower-income location is stated.'' These are different claims with different policy implications, and the present design cannot separate them.

Third, expanding beyond twenty cities would address the most important bound on generalisability. Twenty cities imposes an effective sample size of twenty for every geographic contrast. A design of sixty or more cities, balanced across income tiers and WHO regions, would allow city-income matching --- comparing cities with similar income but different crisis-service infrastructure, and vice versa --- which is the analysis that separates correlation from mechanism.

\section{Concluding Remarks}

Within this sample, the automated analyses indicate geographic variation in actionability, surface localisation and religious framing, while the crisis-contact audit identifies visibly suspicious patterns concentrated in less-served settings. These findings should be read according to their evidential status. RQ1 and RQ2 are exploratory because their measures failed fresh hold-out validation; RQ3 has separate validation support but remains single-coder and geographically confounded. None proves global bias, stereotyping or fabrication.

The methodological contribution may prove as durable as the substantive findings. LLM audit studies that rely on automated content coding without independent human validation face a specific and underappreciated risk: the measure may agree with itself while measuring the wrong thing. This study shows that risk is not hypothetical, and that a validation protocol capable of catching directional failures --- not merely calibrating thresholds --- is both feasible within a dissertation-scale project and necessary for defensible inference.
